%% Materials and experimental setup (MedIA revision, v2 pipeline).
%% Counts are read from the fold-cache manifests
%% ($V2/cache{64,96}/{cohort}/fold*/manifest.json); training details from src/v2/.
\section{Materials and experimental setup}
\label{sec:setup}

\subsection{Cohorts}
\label{sec:setup-cohorts}

The four primary cohorts form a $2\times2$ design that crosses the task (tumor
versus non-tumor tissue, or mitotic figure versus non-mitotic look-alike) with
the shift type. Under an \emph{institutional} shift, domains differ in center,
staining, scanner and patients together. Under an \emph{acquisition-only}
shift, the domains are scanners within one laboratory
(Fig.~\ref{fig:design}, Table~\ref{tab:cohorts}). All splits are row-shuffled with a fixed seed.
Within a fold, the labeled splits are pairwise disjoint in their groups
(patient, slide or image), and the unlabeled pool shares no group with the
validation or test splits. A verification script recomputes these conditions
from the stored group identifiers. It also checks that every labeled split
contains both classes, that ID validation covers every training domain and that
no patch is duplicated within or across splits. Every fold passes.

\begin{figure}[pos=tbp]
  \centering
  \includegraphics[width=\linewidth]{fig_design}
  \caption{Study design. (a) The four cohorts cross the task with the type of
  shift, one cohort per cell. (b) Leave-one-domain-out folds for $K=5$: fold $t$
  tests domain $t$, validates out of distribution on domain $(t-1)\bmod K$ and
  holds out ID-validation groups of every training domain. (c) Every run is
  checkpointed ten times and evaluated under four selection rules; the
  positions of the selected checkpoints are illustrative.}
  \label{fig:design}
\end{figure}

\begin{table}[width=\linewidth,pos=t]
\caption{Cohorts and per-fold split sizes (number of patches; a range is
the minimum--maximum over folds). Counts are identical at 64 and 96~px. ID-val:
group-disjoint validation split of the training domains; OOD-val: held-out
validation domain $(t-1)\bmod K$. The un-normalized MIDOG~2021 variant, used as
a compact-tier sensitivity analysis, has the same split sizes as MIDOG~2021.}
\label{tab:cohorts}
\footnotesize
\begin{tabular*}{\tblwidth}{@{\extracolsep{\fill}}>{\raggedright\arraybackslash}p{2.5cm}*{4}{>{\raggedright\arraybackslash}p{3.05cm}}@{}}
\toprule
 & Camelyon17-WILDS & Canine cSCC & MIDOG 2021 & MIDOG++ \\
\midrule
Task (pos.\ vs neg.) & tumor vs normal (lymph node) & SCC tumor vs non-tumor tissue
  & mitotic figure vs look-alike & mitotic figure vs look-alike \\
Domain & hospital & scanner (same 44 slides) & scanner (one lab) & tumor type \\
$K$ & 5 & 5 & 3 & 7 \\
Shift type & institutional & acquisition-only & acquisition-only
  & institutional and biological \\
Disjoint group & patient & slide & image & image \\
OOD-val domain & yes & yes & none & yes \\
Scale normalization & none (native patch) & per scanner, to CS2 & 0.25\,\textmu m/px & 0.25\,\textmu m/px \\
\midrule
Train & 30,000 & 22,872--26,596 & 2,160--2,544 & 11,854--17,940 \\
ID-val & 12,000 & 3,895--4,780 & 597--716 & 2,555--3,891 \\
OOD-val & 12,000 & 2,197--3,745 & -- & 1,806--8,216 \\
Test & 25,000 & 2,197--3,745 & 1,175--1,648 & 1,806--8,216 \\
Unlabeled pool & 40,000 & 40,000 & 25,200 & 40,000 \\
Positive rate, train & 0.50 & 0.52--0.55 & 0.35--0.39 & 0.42--0.49 \\
Positive rate, test & 0.50 & 0.51--0.59 & 0.35--0.43 & 0.27--0.63 \\
\bottomrule
\end{tabular*}
\end{table}

\paragraph{Cohort construction} Camelyon17-WILDS provides $96\times96$~px
lymph-node patches from five hospitals \citep{koh2021wilds,bandi2019camelyon17};
ID validation holds out whole patients of every training hospital. In the
canine cohort, each of 44 slides was digitized with five scanners
\citep{wilm2023canine}. The slides are split into five groups, so that a test
domain is a new scanner and a set of slides that is absent from training under
every scanner, and magnification is normalized to the CS2 scanner. MIDOG~2021
\citep{aubreville2023midog} has 150 annotated breast cancer images from one
laboratory, 50 per scanner. The scanners' resolutions differ by about 11\%, so
the primary cache normalizes every image to 0.25~\textmu m/px; the
un-normalized cache is a compact-tier sensitivity analysis. MIDOG++
\citep{aubreville2023midogpp} provides 503 images of seven tumor types from two
species, and its 150 breast cancer images are the MIDOG~2021 images. Mitosis
patches are centered on annotated mitotic figures and non-mitotic look-alikes.
The unlabeled pools used for self-supervision come only from the training
domains and exclude all validation and test groups, so the self-supervised arms
remain domain generalization methods. Sampling, scale normalization, the
MIDOG++ domains and the handling of objects near image borders are detailed in
Supplementary Section~\ref{S-sup:setup}.

\subsection{Model tiers}
\label{sec:setup-models}

\paragraph{Compact tier (C)} The compact model is a four-stage convolutional
network on 64~px input with 583,394 parameters. Stages 1--3 each apply two
$3\times3$ convolution--batch normalization--ReLU blocks followed by $2\times2$
max pooling, with 32, 64 and 128 channels. Stage 4 is one such block with 256
channels, followed by global average pooling and a linear layer with two
outputs. It is trained from random initialization, and inputs are normalized
with ImageNet channel statistics.

\paragraph{Standard tier (S)} ResNet-50 \citep{he2016resnet} with ImageNet
initialization \citep{deng2009imagenet} is fine-tuned end to end at the native
96~px. Two pathology foundation models of \citet{kang2023benchmarking} are also
used: a Barlow Twins ResNet-50, fine-tuned at 96~px, and a DINO ViT-S/16
\citep{dosovitskiy2021vit}, fine-tuned on inputs upsampled bilinearly to
224~px. Their inputs are normalized with the channel statistics released with
the weights. Frozen linear probes on ImageNet ResNet-50, Barlow Twins ResNet-50
and DINO ViT-S/16 follow one rule: the input is upsampled to 224~px, and
features are extracted in fp32 and standardized with the training-split mean
and SD. A linear head is trained with SGD (learning rate 0.05, momentum 0.9,
weight decay $10^{-4}$, 1,200 steps of 128) and selected like every other arm.
A supplementary representation ladder applies the same probe to nine frozen
backbones. \citet{kang2023benchmarking} pre-trained on 19 million patches from
20,994 TCGA whole-slide images, and the released checkpoints are these
TCGA-only models. None of our cohorts comes from TCGA. Camelyon17 comes
from Dutch hospitals, the canine slides from the veterinary CATCH collection
\citep{wilm2023canine}, and
MIDOG++ from veterinary laboratories and UMC Utrecht. The pre-training images
therefore do not overlap with ours, although some organs (breast, skin) occur in
both.

\subsection{Optimization and checkpointing}
\label{sec:setup-optim}

Training uses the fixed budget of Section~\ref{sec:protocol-lodo}: batches of
128 drawn without replacement and reshuffled every epoch, and a cosine schedule.
The compact tier uses SGD (momentum 0.9, weight decay $10^{-4}$) with a fixed,
untuned learning rate of 0.02. The standard-tier ResNet-50 arms use SGD with
the same momentum and weight decay. Their learning rate is selected per fold
from $\{10^{-3},3\times10^{-3},10^{-2},3\times10^{-2},10^{-1}\}$ on ImageNet
ERM and shared by every ImageNet ResNet-50 arm of the fold. The grid was extended
upward twice because the selected rate sat on its upper edge. The Barlow Twins
fine-tunes have their own per-fold rate, selected on Barlow Twins ERM from
$\{3\times10^{-4},10^{-3},3\times10^{-3},10^{-2},3\times10^{-2}\}$. Inheriting
the ImageNet rate made their training fail on some folds
(Supplementary Section~\ref{S-app:tuning}). The ViT uses
AdamW (weight decay 0.05, 5\% linear warm-up) with a learning rate selected per
fold from $\{10^{-5},3\times10^{-5},10^{-4}\}$. Every arm except the frozen
probes receives a random dihedral augmentation (flips and rotations by
multiples of $90^\circ$, drawn per batch). The standard tier trains in
bfloat16 autocast and the compact tier in fp32. Evaluation is always in fp32,
with float64 probabilities. TF32 is enabled, and kernels are not forced to be
deterministic. At each checkpoint the record stores eight metrics per split
(AUROC, average precision, accuracy, balanced accuracy, F1, ECE, Brier score,
negative log-likelihood) and the ID-validation AUROC of each training domain.
Predictions at the selected checkpoints and the ID-selected weights are also
saved. A run whose loss or predictions become non-finite is stopped, marked as
diverged and kept in the analysis, with non-finite probabilities scored as 0.5.
Six of the 3,400 final runs diverged, all IRM: two at the compact tier and
four at the standard tier. Camelyon17 models are also scored on the
PatchCamelyon test split \citep{veeling2018pcam}, center-cropped to 64~px at the
compact tier. This score is never used for selection.

\subsection{Arms}
\label{sec:setup-arms}

Table~\ref{tab:arms} lists every arm with its implementation and grid. Arms
that combine two components inherit the tuned value of their parent:
SimCLR+SAM takes $\rho$ from SAM, and the foundation-model H\&E-jitter arms
take $s$ from standard-tier H\&E jitter.

\begin{table}[width=\linewidth,pos=p]
\caption{Arms, implementations and hyperparameter grids. Grids are searched
per fold with seed~0 and selected by ID-validation AUROC; values outside a
grid are fixed. Tier: C = compact CNN (64~px), S = standard (96~px). All
fine-tuned arms also receive the shared dihedral augmentation. $^{\dagger}$Transductive: uses the unlabeled images of the test domain.}
\label{tab:arms}
\footnotesize
\begin{tabular*}{\tblwidth}{@{\extracolsep{\fill}}>{\raggedright\arraybackslash}p{2.4cm}>{\raggedright\arraybackslash}p{7.4cm}>{\raggedright\arraybackslash}p{4.4cm}l@{}}
\toprule
Arm & Implementation & Grid / fixed values & Tier \\
\midrule
\multicolumn{4}{@{}l}{\emph{Baseline and controls}}\\
ERM & Cross-entropy. & C: lr 0.02; S: lr $\in\{10^{-3},3{\times}10^{-3},\dots,10^{-1}\}$ & C, S \\
Augmentation-only & SimCLR's view augmentation (random resized crop, flips, rotations, color jitter, grayscale) applied to supervised batches; no contrastive loss. & crop scale $[0.4,1]$; jitter 0.4 with prob.\ 0.8; hue offset 0.08; grayscale prob.\ 0.2 & C \\
Compute-matched ERM & ERM for $1{,}200+\lceil I_{\text{SSL}}/128\rceil$ steps; $I_{\text{SSL}}$ = images passed through the encoder during SimCLR pre-training (two views per image). & 4,944 steps ($|U|=40{,}000$); 3,552 steps (MIDOG~2021, $|U|=25{,}200$) & C \\
\midrule
\multicolumn{4}{@{}l}{\emph{Optimization and self-supervision}}\\
SAM \citep{foret2021sam} & SGD base optimizer; batch-norm statistics updated on the unperturbed pass only. & $\rho\in\{0.02,0.05,0.1\}$ & C, S \\
SimCLR \citep{chen2020simclr} & NT-Xent pre-training on the fold's unlabeled pool (C: from scratch; S: continued from ImageNet), 2-layer projection head, then ERM fine-tuning. & $\tau=0.5$, batch 256, 6 epochs, lr 0.05 (C) / 0.01 (S) & C, S \\
SimCLR+SAM & SimCLR initialization, SAM fine-tuning. & $\rho$ from SAM & C \\
\midrule
\multicolumn{4}{@{}l}{\emph{Domain-aware training (uses training-domain labels)}}\\
GroupDRO \citep{sagawa2020groupdro} & Weights over training domains, $q_k\leftarrow q_k e^{\eta L_k}$, loss $\sum_k q_kL_k$; domains absent from a batch are not updated. & $\eta\in\{10^{-3},10^{-2},10^{-1}\}$ & C, S \\
IRMv1 \citep{arjovsky2019irm} & Split-batch penalty per domain, averaged over domains; weight 1 for the first 25\% of steps, then $\lambda$, with optimizer reset. & $\lambda\in\{1,10,100\}$ & C, S \\
DeepCORAL \citep{sun2016coral} & Squared differences of feature means and covariances, averaged over pairs of training domains. & $\lambda\in\{0.1,1,10\}$ & C, S \\
MixStyle \citep{zhou2021mixstyle} & Random mixing of instance feature statistics after stages 1 and 2 (C) or ResNet layers 1 and 2 (S); identity at evaluation. & $p\in\{0.5,1\}\times\alpha\in\{0.1,0.3\}$ & C, S \\
Fish \citep{shi2022fish} & Inner copy takes one SGD step per training domain in the batch (random order, persistent inner optimizer); $\theta\leftarrow\theta+\epsilon(\tilde\theta-\theta)$ for parameters, batch-norm buffers copied. & $\epsilon\in\{0.01,0.05,0.1,0.5\}$ & C, S \\
LISA \citep{yao2022lisa} & Per batch, intra-label mixup across domains (prob.\ $p_{\text{sel}}$) or intra-domain mixup across labels; $\lambda\sim\mathrm{Beta}(\alpha,\alpha)$, soft labels. & $\alpha\in\{0.5,2\}\times p_{\text{sel}}\in\{0.5,1\}$ & C, S \\
\midrule
\multicolumn{4}{@{}l}{\emph{Stain normalization and color augmentation}}\\
Macenko$^{\dagger}$ \citep{macenko2009} & One stain estimate per slide and scanner, mapped to a reference fitted with the same estimator on the training slides; every split. & $I_0=240$, OD threshold 0.15, angle percentile 1; $\le$256 tiles per slide & C, S \\
H\&E jitter & Per-image brightness, contrast and saturation factors $\sim U[1-s,1+s]$ and per-channel offsets $\sim U[-s/5,s/5]$. & $s\in\{0.2,0.4\}$ & C, S \\
TIA-style$^{\dagger}$ & Per-slide Macenko normalization followed by H\&E jitter. & $s=0.3$ & C, S \\
Per-tile Macenko & One stain estimate per tile (Supplementary Section~\ref{S-app:macenko}). & as Macenko & C, S \\
\midrule
\multicolumn{4}{@{}l}{\emph{Test-time adaptation (on the ID-selected checkpoint)}}\\
BN-adapt$^{\dagger}$ \citep{schneider2020bn} & Batch-norm statistics re-estimated (cumulative average) on the shuffled test stream in batches of 128, then prediction. & on ERM, SimCLR, H\&E jitter (C); ERM (S) & C, S \\
TENT$^{\dagger}$ \citep{wang2021tent} & Online entropy minimization of batch-norm affine parameters (Adam, lr $10^{-3}$, one step per batch of 128, shuffled stream); a batch is predicted before its update. & on ERM & C \\
\midrule
\multicolumn{4}{@{}l}{\emph{Pathology foundation models \citep{kang2023benchmarking}}}\\
Barlow Twins RN50: ERM, +H\&E jitter & Fine-tuned at 96~px. & lr $\in\{3{\times}10^{-4},\dots,3{\times}10^{-2}\}$ on BT ERM; $s$ from H\&E jitter & S \\
DINO ViT-S/16: ERM, +H\&E jitter & Fine-tuned at 224~px, AdamW. & lr $\in\{10^{-5},3{\times}10^{-5},10^{-4}\}$; $s$ from H\&E jitter & S \\
Frozen probes & Linear head on frozen ImageNet RN50, Barlow Twins RN50 or DINO ViT-S/16 features at 224~px. & lr 0.05 & S \\
\bottomrule
\end{tabular*}
\end{table}

\paragraph{Controls for self-supervision} The augmentation-only control
applies SimCLR's view augmentation to labeled batches without the contrastive
loss, which separates the effect of the augmentations from that of
pre-training. The compute-matched control gives ERM as many images through the
encoder as SimCLR receives in pre-training and fine-tuning together
(Table~\ref{tab:arms}).

\paragraph{Stain normalization} As in common practice, Macenko's stain matrix
\citep{macenko2009} is estimated once per slide imaged on one scanner, from the
pooled tissue pixels of up to 256 of its tiles, with concentrations clamped at
zero, and is mapped to a reference fitted with the same estimator on the
training slides. No labels are used, but the method is transductive: each test
slide is normalized from its own tiles. The per-tile variant estimates the
stain matrix on every tile separately. Implementation details are in
Supplementary Section~\ref{S-sup:setup}, and the effect of the granularity in
Supplementary Section~\ref{S-app:macenko}. Our H\&E jitter perturbs RGB color; it
is not an augmentation in HED space.

\paragraph{Test-time adaptation} BN-adapt and TENT process the test split in a
seed-specific random order. Each batch of 128 therefore contains both classes
in roughly the test domain's proportion. A label-sorted stream would give
single-class batches and an uninformative adapted AUROC.

\paragraph{Transductive arms} BN-adapt, TENT, Macenko and TIA-style use the
unlabeled images of the test domain at prediction time; the other arms see one
test image at a time. The four transductive arms are marked in every table and
figure and are not directly comparable with the inductive arms.

\paragraph{Dropped arm} The rotation-invariance arm of the previous version was
dropped. In this pipeline it reduced to the dihedral augmentation that every arm
receives, so it was identical to ERM.

\subsection{Seeds, compute and availability}
\label{sec:setup-compute}

Final runs use seeds 0--4 at the compact tier and 0--2 at the standard tier.
Tuning uses seed~0 and is stored separately from the final runs. The planned
design has the following runs:
\begin{itemize}
  \item tuning: 520 at the compact tier (26 configurations, 20 folds) plus 78
    for the un-normalized MIDOG~2021 cohort; at the standard tier, 160
    learning-rate runs (100 ResNet-50, 60 ViT), 100 Barlow Twins
    learning-rate runs and 520 method-grid runs;
  \item final: 1,600 at the compact tier (16 arms, 5 seeds, 20 folds), 240 for
    the un-normalized MIDOG~2021 sensitivity analysis, and 1,200 at the standard
    tier (20 arms, 3 seeds, 20 folds);
  \item representation ladder: 360 runs.
\end{itemize}
Test-time adaptation is evaluated within the run of the base arm. All 4,778 runs
(1,378 tuning and 3,400 final, the ladder included) completed. They ran on
NVIDIA H100 NVL GPUs and used about 34 GPU-hours, not counting development
and diagnostic runs. Each
record contains the full hyperparameters, a SHA-256 hash of the training code
and a fingerprint of the fold cache. A record whose fingerprint or
hyperparameters do not match the job is recomputed. Code, fold-construction
scripts, all run records and the scripts that regenerate every table and figure
are available at \coderepo\ \TBD{Zenodo DOI of the archived release}.
